% Cambio de base: las coordenadas son etiquetas (partición por clase: apunte/scripts/reorganizar_por_clase.py)
\chapter{Cambio de base: las coordenadas son etiquetas}\label{cap:cambio-base}

\section{La idea}

\section{Construcciones, cuentas y ejemplos}

\section{En profundidad}

\subsection{Cambio de base y diagonalización}\label{sec:cambio-base}

{\small\vspace{1mm}\linkUnidad{08_cambio_de_base}{Simulación interactiva: Unidad 12 — Cambio de base y diagonalización}\par}
{\small\vspace{1mm}\linkAnimacion{03_cambio_de_base_viaje}{Animación: el viaje \(A\,T_{\mathrm{diag}}\,A^{-1}\)}\par}


El cambio de base y la diagonalización se desarrollan en el párrafo ``La matriz de cambio de
base y el despeje de la diagonalización''; después se extienden a potencias \(A^t\) y trayectorias.

\paragraph{Hecho geométrico: la flecha no se movió, cambió la grilla.} Un mojón tiene
coordenadas \((3,4)\) en la grilla cuadrada de siempre. Si en vez de los ejes canónicos usamos
las direcciones \(b_1=(1,1)\) y \(b_2=(1,-1)\), el \emph{mismo} punto se escribe
\[
  \begin{pmatrix}3\\4\end{pmatrix}
  =\tfrac72\begin{pmatrix}1\\1\end{pmatrix}-\tfrac12\begin{pmatrix}1\\-1\end{pmatrix},
  \qquad\text{coordenadas nuevas }(\tfrac72,-\tfrac12).
\]
El mojón está donde estaba; cambió la regla con que lo medimos. Escalera: en \(\R\) cambiar
de base es cambiar de unidad (de metros a centímetros el número se multiplica por 100 y la
longitud no cambia); en \(\R^2\) la grilla puede ser oblicua y con pasos distintos; en
\(\R^n\) es una matriz \(n\times n\) invertible. No parece
haber un conjunto de vectores privilegiado a priori: lo único que distingue a la base
canónica es que con ella el producto interno es la cuenta de siempre.

\begin{resultado}[La matriz de cambio de base]
Si \(P\) tiene en sus columnas los vectores de la base nueva escritos en la canónica,
\[
  [v]_{\text{can}}=P\,[v]_{\text{nueva}},\qquad [v]_{\text{nueva}}=P^{-1}[v]_{\text{can}} .
\]
Con \(P=\begin{pmatrix}1&1\\1&-1\end{pmatrix}\), \(P^{-1}=\tfrac12\begin{pmatrix}1&1\\1&-1\end{pmatrix}\)
y \(P^{-1}(3,4)^T=(\tfrac72,-\tfrac12)^T\), como arriba. Los tres tipos de objeto cambian
de base distinto, y conviene escribir siempre el sentido (acá, de la canónica a la nueva):
los vectores como \([v]_{\text{nueva}}=P^{-1}[v]_{\text{can}}\); las formas lineales, que
son filas, multiplicando por \(P\) del lado derecho, \([\varphi]_{\text{nueva}}=[\varphi]_{\text{can}}\,P\)
(sección~\ref{sec:formas-lineales}), para que \(\varphi(v)\) dé lo mismo en las dos bases;
y las transformaciones con
\[
  A_{\text{nueva}}=P^{-1}\,A\,P ,
\]
que se lee de derecha a izquierda: pasamos las coordenadas nuevas a canónicas, aplicamos \(A\)
como siempre, volvemos a las nuevas. Todo se convierte a la vez para que ninguna relación
cambie: cada punto sigue en su recta y \(A\) sigue llevando cada punto adonde lo llevaba.
\end{resultado}

\paragraph{Ejercicio resuelto: grilla de formas y vector.} Consigna: formas lineales \((1,1)\) y \((-1,2)\), vector \((3,4)\), transformación
\(T=\begin{pmatrix}1{,}2&0{,}5\\-0{,}1&1{,}3\end{pmatrix}\), elegida no muy distinta de
la identidad. Dibujar las curvas de nivel de las dos formas (dos mazos de rectas
paralelas que juntos arman una grilla), ubicar el vector en la grilla, transformar el
vector con \(T\) y las formas con \(T^{-1}\), redibujar y verificar que el vector
transformado cae en la misma intersección. Antes de transformar, \((1,1)\cdot(3,4)=7\) y
\((-1,2)\cdot(3,4)=5\): el vector está en la séptima curva de una y la quinta de la otra.
Inversa por adjunta sobre determinante, \(\det T=1{,}56+0{,}05=1{,}61\):
\[
  T^{-1}=\frac1{1{,}61}\begin{pmatrix}1{,}3&-0{,}5\\0{,}1&1{,}2\end{pmatrix},\qquad
  T\begin{pmatrix}3\\4\end{pmatrix}=\begin{pmatrix}5{,}6\\4{,}9\end{pmatrix},
\]
\[
  (1,1)\,T^{-1}=\tfrac1{1{,}61}(1{,}4,\ 0{,}7)=(0{,}870,\ 0{,}435),\qquad
  (-1,2)\,T^{-1}=\tfrac1{1{,}61}(-1{,}1,\ 2{,}9)=(-0{,}683,\ 1{,}801).
\]
Verificación: \(0{,}870\cdot5{,}6+0{,}435\cdot4{,}9=7{,}0\) y
\(-0{,}683\cdot5{,}6+1{,}801\cdot4{,}9=5{,}0\). El vector transformado sigue en la séptima
curva de la primera forma y en la quinta de la segunda. La forma \((1,1)\)
queda \((0{,}86,\ 0{,}43)\), se acuesta un poco (redondeo), y la \((-1,2)\) queda un poco más
vertical. Hecho geométrico: la grilla se deformó con el plano y el punto se movió con el
plano; medido con la grilla deformada, el punto está donde estaba. No hacemos nada: convertimos
todo para que nada cambie. En planilla: dos columnas
para la forma, dos para el vector, \texttt{SUMAPRODUCTO}.

\paragraph{Diagonalizar es elegir la grilla de los autovectores.} Sea
\(A=\begin{pmatrix}2&1\\1&2\end{pmatrix}\). Polinomio característico
\((2-\lambda)^2-1=0\), autovalores \(3\) y \(1\), autovectores \((1,1)\) y \((1,-1)\)
(perpendiculares porque \(A\) es simétrica). Con esos autovectores como columnas de \(P\),
\[
  P^{-1}AP=\tfrac12\begin{pmatrix}1&1\\1&-1\end{pmatrix}
  \begin{pmatrix}3&1\\3&-1\end{pmatrix}
  =\begin{pmatrix}3&0\\0&1\end{pmatrix}=D .
\]
Hecho geométrico: en la grilla de los autovectores la transformación es dos estiramientos
independientes, por \(3\) a lo largo de \(b_1\) y por \(1\) a lo largo de \(b_2\); la grilla
``se estira o se comprime pero no se rompe''. En la grilla común la misma transformación
no tiene lectura fácil. Equivale a aplicar dos afinidades
axiales sucesivas.

\paragraph{La matriz de cambio de base y el despeje de la diagonalización.} En este
párrafo usamos otros nombres: \(A\) es la matriz que tiene los autovectores en
columnas (la \(P\) del recuadro de arriba, con otro nombre) y \(T\) es la transformación
(la \(A\) del párrafo anterior).

\emph{El diagrama de conmutatividad.} Dualidad y cambio de base se juntan en una sola
cadena, escrita en orden de aplicación (de izquierda a derecha):
\[
  (p)\;\to\;(A)\;\to\;[\,A^{-1}\to T_{\mathrm{diag}}\to A\,]\;\to\;(A^{-1})\;\to\;(\ell) .
\]
En un extremo entra un punto \(p\) y en el otro una recta \(\ell\), los dos escritos en la
base de autovectores; lo que se mide es la incidencia. Con la convención de abajo (\(A\)
lleva de la base de autovectores a la canónica), el punto pasa a la canónica como \(A\,p\)
y la forma como \(\ell\,A^{-1}\) (multiplicando del lado derecho). El recuadro del medio
recibe y devuelve coordenadas canónicas: adentro actúa \(T_{\mathrm{diag}}\), en la base
de autovectores, y el recuadro entero, leído como producto (lo que se aplica último va a la
izquierda), es \(A\,T_{\mathrm{diag}}\,A^{-1}=T_{\mathrm{can}}\), la misma transformación
escrita en la base canónica. La cadena completa es
\(\ell A^{-1}\,(A\,T_{\mathrm{diag}}\,A^{-1})\,A\,p=\ell\,T_{\mathrm{diag}}\,p\): las
conversiones se absorben de a pares (\(A^{-1}A=I\)) y la incidencia no cambia; el punto
transformado sigue cayendo en la recta transformada. Despejar
\(T_{\mathrm{diag}}=A^{-1}\,T_{\mathrm{can}}\,A\) es sacar la diagonal del recuadro, como
hacemos más abajo. Todo se convierte a la vez para que nada cambie.

\emph{La matriz con los autovectores en columnas.} Definimos \(A=[\,v_1\mid v_2\,]\), los
autovectores escritos en la base canónica, uno por columna, y la probamos con el vector
\((1,0)^T\):
\[
  A\begin{pmatrix}1\\0\end{pmatrix}
  =\begin{pmatrix}A_{11}&A_{12}\\A_{21}&A_{22}\end{pmatrix}\begin{pmatrix}1\\0\end{pmatrix}
  =\begin{pmatrix}A_{11}\\A_{21}\end{pmatrix}=v_1 .
\]
Entra una unidad del primer autovector escrita en la base de autovectores y sale ese mismo
autovector escrito en la canónica: \(A\) traduce de la base de autovectores a la canónica
y \(A^{-1}\) hace el camino inverso, de la canónica a la de autovectores. El recorrido
completo para un punto cualquiera: el punto \((1,1)=e_x+e_y\), dado en la canónica, pasa
con \(A^{-1}\) a coordenadas \((c_1,c_2)\) en la base de autovectores, es decir
\((1,1)=c_1v_1+c_2v_2\); ahí la transformación es diagonal y devuelve
\((\lambda_1c_1,\ \lambda_2c_2)\); y \(A\) traduce ese resultado de vuelta a la canónica,
\(\lambda_1c_1v_1+\lambda_2c_2v_2\), que es \(T(1,1)\). Con la matriz del párrafo
anterior, \(v_1=(1,1)\) y \(v_2=(1,-1)\), el punto \((1,1)\) es el propio \(v_1\):
\((c_1,c_2)=(1,0)\), la diagonal da \((3,0)\) y \(A\) devuelve \(3v_1=(3,3)\), que
coincide con \(\begin{pmatrix}2&1\\1&2\end{pmatrix}(1,1)^T=(3,3)^T\).

\emph{El despeje.} El mismo recorrido, escrito como producto de matrices, dice que la
transformación en la base canónica es la diagonal conjugada por el cambio de base:
\[
  T_{\mathrm{can}}=A\,T_{\mathrm{diag}}\,A^{-1}
  \quad\Longrightarrow\quad
  T_{\mathrm{diag}}=A^{-1}\,T_{\mathrm{can}}\,A ,
\]
leída de derecha a izquierda: \(A^{-1}\) lleva a la base de autovectores,
\(T_{\mathrm{diag}}\) estira cada eje propio, \(A\) vuelve a la canónica. El despeje se
hace multiplicando por \(A^{-1}\) a la izquierda y por \(A\) a la derecha; los dos
productos \(A^{-1}A=I\) y \(A\,A^{-1}=I\) que aparecen se cancelan y queda
\(T_{\mathrm{diag}}\) sola. Con los nombres de este apunte (\(P\) para la matriz de
autovectores, \(A\) para la transformación) es exactamente el \(D=P^{-1}AP\) del párrafo
anterior: misma matriz, otro nombre.

\paragraph{Potencias y la trayectoria de un punto.} Volvemos a los nombres del apunte:
\(P\) es la matriz de autovectores y \(A\) la transformación. Como \(A=PDP^{-1}\),
\(A^n=PD^nP^{-1}\): aplicar \(n\) veces la transformación es estirar \(n\) veces cada eje
propio. Para \(n=5\),
\[
  A^5=P\begin{pmatrix}243&0\\0&1\end{pmatrix}P^{-1}
  =\tfrac12\begin{pmatrix}244&242\\242&244\end{pmatrix}
  =\begin{pmatrix}122&121\\121&122\end{pmatrix}.
\]
Lectura: \((1,0)=\tfrac12(1,1)+\tfrac12(1,-1)\); la parte en \((1,1)\) se multiplica por
\(243\) y la parte en \((1,-1)\) queda igual: \(\tfrac12(243,243)+\tfrac12(1,-1)=(122,121)\),
la primera columna. Si los autovalores son positivos el exponente puede ser real:
\(A^t=PD^tP^{-1}\) con \(D^t=\operatorname{diag}(3^t,1)\), y el punto \((1,0)\) recorre
\[
  A^t\begin{pmatrix}1\\0\end{pmatrix}=\tfrac12\begin{pmatrix}3^t+1\\3^t-1\end{pmatrix}:
  \quad(1,0)\ \text{en } t=0,\quad(1{,}37,\,0{,}37)\ \text{en } t=\tfrac12,\quad(2,1)\ \text{en } t=1 .
\]
Es la tercera idea del hilo (sección~\ref{sec:hilo}): la transformación se identifica con
la trayectoria que genera, y la potencia real la convierte en película. La misma jugada que
la exponencial de una rotación, ahora con estiramientos en vez de giros. La unidad
\texttt{08\_cambio\_de\_base} (apéndice~\ref{ap:notebooks}) tiene el deslizador en \(t\).

\paragraph{Dirección de cada cambio de base.} En la factorización \(A=P\,D\,P^{-1}\), el factor de la derecha, el que actúa primero sobre un vector en canónica, es el que \emph{convierte de la canónica a la base de autovectores}. Con subíndices descriptivos:
\(A\,T_{\mathrm{diag}\langle Av\rangle}A^{-1}=T_{\mathrm{no\,diag}\langle\mathrm{can}\rangle}\)
con \(A\) = autovectores en columnas, que es \(P\,D\,P^{-1}\) de arriba. Para no mezclar, en este apunte \(P\) siempre tiene la base nueva en columnas y se lee ``\(P\) lleva de la nueva a la canónica''.

\begin{resultado}[Diagonalizabilidad no tiene nada que ver con invertibilidad: cuatro matrices testigo]
Las dos preguntas, ``¿qué pasa si hay algo en el núcleo?'' y ``¿qué pasa si
un autovalor repetido no alcanza para armar una base?'', son independientes. Invertir
pide que el núcleo sea \(\{0\}\), es decir, que ningún autovalor sea \(0\). Diagonalizar
pide una base de \(\R^2\) formada por autovectores. Las cuatro combinaciones aparecen:
\medskip

\renewcommand{\arraystretch}{1.35}
\noindent\begin{tabular}{@{}p{2.6cm}p{1.1cm}p{2.2cm}p{2.1cm}p{5.6cm}@{}}
\textbf{Matriz} & \(\det\) & \textbf{Autovalores} & \textbf{¿Base de autovectores?} & \textbf{Hecho geométrico}\\
\hline
\textbf{1.}~\(\bigl(\begin{smallmatrix}2&0\\0&1\end{smallmatrix}\bigr)\) & \(2\ne0\) & \(2,\ 1\) & sí, \(e_1,e_2\) &
diagonalizable e invertible: estira el eje \(x\) al doble y deja fijo el \(y\); dos direcciones propias independientes.\\
\textbf{2.}~\(\bigl(\begin{smallmatrix}1&0\\0&0\end{smallmatrix}\bigr)\) & \(0\) & \(\lambda_1=1,\ \lambda_2=0\) & sí, \(e_1,e_2\) &
proyección sobre el eje \(x\): aplasta el eje \(y\) al origen, sin inversa, pero con base completa de autovectores.\\
\textbf{3a.}~\(\bigl(\begin{smallmatrix}1&1\\0&1\end{smallmatrix}\bigr)\) & \(1\) & \(1\) (doble) & no: sólo \(e_1\) &
corte: invertible, con inversa \(\bigl(\begin{smallmatrix}1&-1\\0&1\end{smallmatrix}\bigr)\); sólo el eje \(x\) queda fijo.\\
\textbf{3b.}~\(\bigl(\begin{smallmatrix}0&-1\\1&0\end{smallmatrix}\bigr)\) & \(1\) & \(\pm i\) & no en \(\R\) &
giro de \(90^\circ\): invertible, ninguna recta real queda fija.\\
\textbf{4.}~\(\bigl(\begin{smallmatrix}0&1\\0&0\end{smallmatrix}\bigr)\) & \(0\) & \(0\) (doble) & no: sólo \(e_1\) &
ni diagonalizable ni invertible: manda \(e_2\mapsto e_1\) y \(e_1\mapsto0\).\\
\end{tabular}
\renewcommand{\arraystretch}{1.18}
\medskip

\textbf{Matriz 2.} Tener autovalor \(0\) impide invertir (el núcleo es el eje \(y\),
\(\det=0\)) y no impide diagonalizar: \(e_1\) y \(e_2\) son autovectores con
\(\lambda_1=1\), \(\lambda_2=0\), y en esa base la matriz ya es diagonal.

\textbf{Matriz 3a.} El polinomio característico es \((\lambda-1)^2=0\): la raíz
\(\lambda=1\) tiene multiplicidad algebraica \(2\). Pero \((A-I)v=0\) da \(v_2=0\), de modo
que el autoespacio es sólo el eje \(x\): multiplicidad geométrica \(1\). Falta un segundo
autovector, y sin él no hay base.

\textbf{Matriz 3b.} En el giro el motivo es otro: no hay rectas reales fijas, los
autovalores \(\pm i\) no son reales.

\textbf{Matriz 4.} \(\lambda=0\) doble, pero el núcleo es sólo el eje \(x\): un único
autovector \(e_1\).

El defecto de las matrices 3a, 3b y 4 es exactamente por lo que la sección~\ref{sec:svd}
pasa a la SVD, \(A=U\Sigma V^T\): con una base ortonormal a la entrada (\(V\)) y otra
\emph{distinta} a la salida (\(U\)), toda matriz, incluso un corte o una rectangular, se
lee como diagonal. Es diagonal entre dos grillas distintas, no una semejanza: \(\Sigma\) no
es ``\(A\) en otra base'', los \(\sigma\) no son autovalores (en el corte,
\(\sigma=(\sqrt5\pm1)/2\) y \(\lambda=1\) doble) y no sirve para potencias,
\(A^n\neq U\Sigma^nV^T\). Las dos descomposiciones coinciden sólo si \(A\) es simétrica
semidefinida positiva.
\linkUnidad{09_cambio_de_base_svd}{Simulación interactiva: cambio de base y SVD}
\end{resultado}
